% !TEX root = ../main.tex
% (this is to make TexShop know where your parent file is)

% Context
% Clear evidence of expansive research and deep critical reflection that has informed the creation of the audio to fit within the stylistic norms of the genre or industry sector in question.The work perfectly meets the brief.


In the section where the knight is first revealed and slowly stands up, I've used a chord progression that ends with two major chords separated by a minor third [00:57--01:00]. This evokes a "superhero" kind of feeling due to it's association with the DC and Marvel universes \autocite[04:22-04:48]{8-bitmusictheorySpiderManEffectiveUse2018}. \footnote{This is also the same interval used in the Rebel Fanfare from the Star Wars movies \autocite[01:42--01:45]{flashmusicStarWarsRebel2018}, so it's been around for a while.}


The uplifting feeling then fades into a bittersweet chord progression as the blacksmith invites the knight to strike. The music is somber, yet kinda hopeful, intended to invite the listener to wonder whether this is the blacksmith accepting his end or something else.

I wanted to keep the tension between the sword plunge [01:17] and the breakage [01:26], so I left that part relatively untouched.

The sword shattering boom sounds to me as if it's a low B-ish, so I've adjusted the key of the song to make that fit with the harmony of the chord. Of course, it's not a clear or particularly pitched sound, but it felt a little smoother that way.

Lyrically (see \hyperref[appendixA]{Appendix~A}), I wanted the piece to feel ambiguous, but it's told from the blacksmith's point of view. I've understood this as if the knight is far from the first one who's come knocking on the blacksmith's proverbial door. Judging from the arrows displayed in the village, this is a semi-regular occurrence, and I think the blacksmith is kinda over it. As if the real people who have an issue with the blacksmith aren't the ones who come face to face with him. Instead, they send mercenaries and adventurers, none of whom really have the full context of what is going on. In a way, they're pawns in a game they don't even know that they're playing, hence the "puppet on a string" and "only you can see it through" lines.

The lyrics also contain a quote attributed to Julius Caesar after crossing the Rubicon: "the die is cast"\autocite{AleaIactaEst2026}, indicating that a point of no return has been reached and that what comes next is uncertain. In the next line, the meaning of that phrase is slightly changed by the words "in iron and blood". This changes the meaning of "cast" from the act of throwing a die to the act of \textit{casting} metal. This is another historical reference to Otto von Bismarck, the \textit{Iron Chancellor}\autocite{OttoBismarck2026}, who said that the great questions of the day would be decided through iron and blood, not diplomacy.\footnote{Incidentally, it's also befitting of the blacksmith being, well, a blacksmith.}
